\scanpage{272}
\uline{Questions.} \deleted{\unsure{?}} Soit $A$ un anneau local \deleted{\unsure{et}} noethérien. On suppose $A$ \uline{hensélien}, et un ``bon anneau'', i.e. $\unsure{\operatorname{Sp}}\widehat A-\unsure{\operatorname{Sp}}A$ est \deleted{\unsure{est}} un \unsure{mor!} régulier.\inserted{p. ex \deleted{\unsure{C'est}} \deleted{\unsure{un}} anneau hensélisé d'un anneau local de la géom. alg. sur $\mathbb C$} Soit $X$ propre sur $A$, et $\widehat X=X\otimes_A\widehat A$.

(i) Est-il vrai que \unsure{tte image universelle} de \unsure{--}, $\operatorname{Pic}(X_{\unsure{n}})$ \deleted{\unsure{?}} [\unsure{on a} $=:\ \unsure{\Longleftrightarrow}$ \unsure{élément} dans $\operatorname{Im}(\operatorname{Pic}(\widehat X)\mathbin{\unsure{-|\to}}\operatorname{Pic}(X_{\unsure{n}}))$] est dans $\operatorname{Im}(\operatorname{Pic}(X)\mathbin{\unsure{-|\to}}\operatorname{Pic}(X_{\unsure{n}}))$ ? En particulier, si $\widehat X/\widehat Y$ est projectif, est-il vrai que $X/Y$ est projectif ?

(ii) Supposons que $\operatorname{supp}R^1f_*(\mathcal O_X)\subset\{y\}$, \deleted{\illegible{1}} p. ex. $X|(Y-\{y\})\simeq Y-\{y\}$. Est-il vrai alors que
\[
\operatorname{Pic}(X)\xrightarrow{\sim}\operatorname{Pic}(\unsure{\widehat Y|\widehat X})\quad ??
\]
\deleted{\unsure{N.B.}}

(iii) Supposons que $Y$ soit \unsure{pt} singulier isolé en $y$. Est-il vrai que
\[
\operatorname{Pic}(Y-\{y\})\xrightarrow{\sim}\operatorname{Pic}(\widehat Y-\{\widehat y\})
\]
[Si (ii) est vrai, et si on peut désingulariser, alors (iii) est vrai\ldots.]. \unsure{Que se passe-t-il si} $y$ n'est pas isolé en $\operatorname{Sing}(Y)$, dans le cas où $A$ est un \unsure{Alg.} analytique ?
\margin{What does \unsure{Mumford's} example in ``Publications IHES'' p. 16 give ??}
\figure{\begin{tikzpicture}[x=1cm,y=1cm]
\draw[->] (2,0)..controls(-1,3)and(-.5,6)..(4.5,6.7);
\draw (2,0)..controls(1.5,-.1)and(1.5,-1)..(2.2,-1.1)--(6,-1.1)..controls(6.5,-1.1)and(6.3,-.4)..(6.8,-.4)--(10,-.4)..controls(10.6,-.3)and(10.8,.8)..(10.4,1.2)--(6.4,1.2)..controls(6.1,1.2)and(6.5,.5)..(6,.3)--(2.3,.3)..controls(2,.3)and(2,.1)..(2,0);
\draw (.4,-1.5)..controls(.6,-2.7)and(.4,-2.7)..(0,-2.9)..controls(.5,-3.1)and(.3,-4.4)..(.6,-4.5);
\end{tikzpicture}}
\mspage{271}
